\section{AlignmentGraph-ISD}
\label{sec:method}

Instructional systems design is a set of mutually constrained artifacts rather than
one long document. ADDIE supplies the process and output contract
~\cite{branch2009addie}; constructive alignment requires objectives, activities,
and assessment to agree~\cite{biggs1996enhancing,biggs2011teaching}. The design
problem is therefore to produce a complete ADDIE program from an underspecified
brief while preserving cross-artifact relations that remain checkable under a
small model and a bounded call budget. AlignmentGraph-ISD addresses that problem by
making those relations shared workflow state.

\subsection{Architecture and Alignment State}

The harness is a LangGraph meta-agent with a fixed \texttt{StateGraph}. Six
specialists handle needs and context, objectives, assessments, activities and
materials, implementation, and evaluation in order. Each receives the brief and
relevant state, writes only its assigned artifacts, and is followed by pedagogical
verification. An aggregator renders the completed state into the benchmark's ADDIE
schema (\Cref{fig:architecture}). The model proposes content; the harness controls
decomposition, state transitions, verification, and assembly.

\begin{figure*}
  \centering
  \includegraphics[width=0.92\textwidth]{figures/architecture.pdf}
  \Description{A LangGraph meta-agent routes a design brief through six specialist
  agents, pedagogical verification, a shared alignment graph, and an ADDIE
  aggregator.}
  \caption{AlignmentGraph-ISD architecture. Specialists coordinate through a shared
  alignment graph; a run returns the ADDIE program and a separate JSON and DOT graph
  dump for educator audit and override.}
  \label{fig:architecture}
\end{figure*}

After step $t$, the graph $G_t=(V_t,E_t)$ contains typed objective, activity,
material, assessment, implementation, and evaluation nodes. Section nodes retain
complete phase artifacts, while element nodes expose individual relations. Edges
record source, target, evidence, rationale, establishment mechanism, and a typed
relation. For an objective $o$, the three central relations define structural
coverage:
\begin{multline}
\label{eq:aligned}
\mathrm{aligned}(o) \iff
\big(\exists a\colon (a,o)\in E_{\mathrm{measures}}\big)\\
\wedge\big(\exists x\colon (x,o)\in E_{\mathrm{supports}}\big)
\wedge\big(\exists e\colon (e,o)\in E_{\mathrm{evaluates}}\big).
\end{multline}
Here $a$, $x$, and $e$ denote an assessment item, learning activity, and evaluation
item. Edge presence makes coverage traceable; semantic validity also requires
content and cognitive-demand compatibility. The graph is thus an auditable claim
about relations, not proof that every recorded relation is sound.

\subsection{Deployment Envelope}
\label{sec:assumptions}

The deployment target is a single 12--16\,GB GPU. At two bytes per parameter, bf16
weights require approximately 1.6\,GB for 0.8B, 4\,GB for 2B, 8\,GB for 4B, and
18\,GB for 9B. The 9B model therefore does not fit, making 4B the largest
in-target checkpoint. The remaining VRAM must hold the KV cache and runtime
overhead, leaving roughly 8k tokens per call. This checkable chain defines the
serving condition independently of experiment timing; 9B is an out-of-target
reference, not a proposed deployment configuration. We assume a supplied brief, a
stateless backbone across calls, and the fixed ADDIE output schema; the harness,
not the backbone, preserves state between calls.

\subsection{Constrained Decomposition and Graph-Context Injection}
\label{sec:decomposition}

Each specialist owns a bounded set of node and edge types. The assessment
specialist, for example, may add assessment items and \emph{measures} edges but
cannot rewrite objectives. Before generation, it receives only the brief material,
stable ISD guidance, and a relation-filtered one-hop graph slice needed for its
work. The activity specialist analogously sees objectives, content modules,
strategy, and media selections. The implementation resolves declared identifiers
directly; it performs neither recursive expansion nor graph traversal. The
$-$GraphContext ablation removes these reads while retaining graph writes.

Mandatory links arise from this constrained specialist contract, not from the mere
existence of a graph. Each downstream item must name its upstream objective; the
harness resolves that identifier and records the typed edge. Consequently, a small
model need not reconstruct the full evolving program from prose in every call, and
a downstream specialist cannot silently alter an upstream premise.

\subsection{Verification, Aggregation, and Audit}
\label{sec:qc}

Structural links alone do not establish pedagogical fit. After each update, the
verifier checks relation coverage and resolvability, Bloom action and declared-level
consistency~\cite{anderson2001taxonomy}, assessment formats and scoring criteria,
active practice, delivery feasibility, implementation guidance, and formative
evaluation. Violations return to the specialist that owns the affected artifact for
bounded, targeted repair. Deterministic fallbacks handle unusable responses and
missing structural references; unresolved warnings remain visible when a safe
local repair is unavailable.

The fixed workflow repeats scoped generation, verification, and repair across all
six concerns, then aggregates the graph into the ADDIE output and runs a residual
audit without repair. A separate JSON dump retains every node, edge, evidence,
rationale, check, and repair, and a DOT rendering makes the settled state visible.
Unlike the execution trajectory, this artifact lets an educator trace which
assessment measures an objective, inspect unsupported relations, and edit or
override the design. It supports professional review rather than replacing it.

\subsection{Tested Mechanisms and Inherited Premise}

The ablation isolates graph-context injection and the verifier, which together
constitute the tested alignment substrate. It does not isolate multi-agent
decomposition: every arm uses the same specialist schedule. The tractability of
narrower subtasks is inherited from decomposed prompting and small-model
orchestration~\cite{zhou2023leasttomost,khot2023decomposed,ferraz2024decrim,
shen2024multillmagent,belcak2025slm}; a monolithic comparison is future work
(\Cref{sec:future-work}). The closest benchmark agents---ReAct-ISD, ADDIE-Agent,
and Dick-Carey-Agent---coordinate cross-artifact information mainly through prompts,
prose, or trajectory state~\cite{isdagentbench,yao2023react,branch2009addie,
dick1978systematic}. AlignmentGraph-ISD instead preserves typed relations as state
that can be queried, repaired, dumped, and overridden.

\subsection{Constructive-Alignment Evaluation Protocol}
\label{sec:protocol}

The second contribution measures whether rendered outputs exhibit signals
corresponding to the three legs of \Cref{eq:aligned}. It extracts objective,
assessment, activity, and evaluation text from every system's saved ADDIE output;
self-declared graph identifiers are used only as validity evidence. The protocol
reports a pre-specified, flat seven-signal panel with no primary endpoint, because
constructive alignment has multiple facets that should not be collapsed or ranked.

Two mechanically independent instrument families carry the panel. Family A uses an
encoder to measure textual correspondence. For source texts $X$ and target texts
$Y$, it computes
\[
  |X|^{-1}\sum_{x\in X}\max_{y\in Y}\max(0,\cos(e_x,e_y)),
\]
clipped above at one. This mean--max \emph{rectified} cosine is a scale choice on
$[0,1]$, not a matching threshold: negative association is placed at the same floor
as orthogonality, and no constant decides whether two texts match. The four signals
are \texttt{objective\_assessment\_similarity},
\texttt{objective\_activity\_similarity},
\texttt{objective\_evaluation\_similarity}, and the reverse-direction
\texttt{assessment\_objective\_similarity}. The first three map respectively to
\emph{measures}, \emph{supports}, and \emph{evaluates}. The reverse signal is
non-directional: a high value may reflect few orphan assessments or mere repetition,
so it is reported but excluded from directional counts.

Family B measures cognitive demand without cosine. It reports
\texttt{objective\_cognitive\_congruence}, the share of objectives whose best
assessment reaches their Bloom level; \texttt{porter\_mean}, agreement between
topic--Bloom distributions; and \texttt{webb\_bloom\_consistency}, item-side
cognitive consistency~\cite{porter2002measuring,fulmer2011estimating,
webb1997criteria}. Topic assignment and best-item selection use argmax, not a
threshold.

Signal names describe the operation computed: the word ``alignment'' names the
protocol and the predicate in \Cref{eq:aligned}, never an individual metric, and no
signal is called ``precision.'' Measurement invokes no LLM, prompt,
sampling step, or matching threshold; once encoder and Bloom versions are pinned,
the scores are fixed functions of the extracted text. We therefore claim measurable
signals of constructive alignment, supported initially by established criteria and
cross-family evidence. The protocol does not by itself validate the construct; an
expert-human tier is required to complete that bridge.
